\chapter{Merger Arbitrage}\label{ch:s1:merger-arbitrage}

A company agrees to be bought for \$40 a share in cash, and its stock trades at \$38.60. The \$1.40 spread, 3.6\%, is what the market charges for the risk
that the deal breaks, in which case the stock falls back toward the \$30 it traded at before. Buying at \$38.60 is selling insurance against that break:
the premium is collected in small amounts deal after deal, and the claims arrive together, because deals break more often when markets fall. Mitchell and
Pulvino measured it over 4\,750 mergers: \omterm{def:s1:merger-arbitrage:arb}{merger arbitrage} returns are uncorrelated with the market in flat and rising markets and correlated with it in
severely falling ones, like the returns from selling index puts, and still earned excess returns of about 4\% a year after costs, once that nonlinearity is controlled for. This chapter prices
\omterm{def:s1:merger-arbitrage:spread}{deal spreads}, builds a portfolio of synthetic deals whose breaks cluster in downturns, and measures what the insurance pays. The build is
\texttt{firm.mergerarb}.

\section{Deal spreads and their pricing}

\begin{definition}[Merger arbitrage, tender offer]\label{def:s1:merger-arbitrage:arb}
\emph{Merger arbitrage}\index{merger arbitrage} buys the shares of a company that has agreed to be acquired, below the offered consideration, and (in a
stock deal) sells short the acquirer's shares it will receive, to earn the difference when the deal completes. A \emph{tender offer}\index{tender offer} is
an offer made directly to a company's shareholders to buy their shares at a stated price, as opposed to a merger voted by the board and shareholders.
\end{definition}

\begin{definition}[Deal spread, break price, implied deal probability]\label{def:s1:merger-arbitrage:spread}
The \emph{deal spread}\index{deal spread} is the offered consideration divided by the target's price, minus one. The \emph{break price}\index{break price}
is the price at which the target would trade if the deal failed. The \emph{implied deal probability}\index{implied deal probability} is the completion
probability that makes the target's price the discounted average of the consideration and the break price: $p = (A - B)/(K - B)$ for price $A$, offer $K$
and break price $B$, ignoring discounting.
\end{definition}

For the \$40 deal, $p = (38.60 - 30)/(40 - 30) = 0.86$: the market prices a 14\% chance of a break. At that probability the trade has no expected profit
($0.86 \times 1.40 = 0.14 \times 8.60$); the arbitrageur earns only if deals break less often than the spread implies. Spreads are quoted annualised because
deals take months: a 3.6\% spread over 84 trading days is 11.3\% a year. The synthetic deals (\texttt{firm.mergerarb.simulate\_deals}) are cash deals
announced at about fifty a year over the synthetic market's ten years, at premiums of 20\% to 40\%, resolving after a median of 110 trading days; their
\omterm{def:s1:merger-arbitrage:spread}{break price} moves with the market; their spread prices a 12\% break probability; their true break probability is 6\% plus 1.5 times the market's fall over
the deal's life. The average spread at entry is 4.73\%, 11.8\% annualised, and the implied completion probability 88\%.

\begin{dated}{2026-09}{US premerger waiting periods}\label{dat:s1:merger-arbitrage:hsr}
Under the Hart-Scott-Rodino premerger notification program, once the parties' filings are complete they must wait 30 days (15 days for a cash \omterm{def:s1:merger-arbitrage:arb}{tender offer}
or a bankruptcy) before consummating the deal, unless the agencies grant early termination; after a preliminary review the reviewing agency (the FTC or the
Department of Justice) can let the period expire, end it early, or issue a Second Request for more information, which extends the review.
\end{dated}

\section{Break risk and its correlation with the market}

\begin{center}
\small
\begin{tabular}{@{}lcccc@{}}
\toprule
synthetic deal portfolio, years 3 to 10 & base case & breaks & spread priced & both\\
 & & independent & at 6\% & \\
\midrule
return a year; over the 4\% rate & 6.76\%; 2.76\% & ---; 4.30\% & ---; $-0.47\%$ & ---; 1.03\%\\
volatility; Sharpe ratio (excess) & 3.26\%; 0.85 & ---; 1.50 & ---; $-0.17$ & ---; 0.46\\
beta in months the market fell over 4\% & 0.20 & 0.20 & 0.17 & 0.17\\
beta in the other months & 0.09 & 0.11 & 0.04 & 0.06\\
completed deals & 92.5\% & 94.7\% & 92.5\% & 94.7\%\\
\bottomrule
\end{tabular}
\end{center}

The base-case portfolio holds 22 deals on average, 374 over the eight years, equally weighted. It earns 2.76\% a year over the rate its spreads discount at,
with a volatility of 3.26\% and a Sharpe ratio of 0.85. Deals completed 92.5\% of the time, against the 88\% the spreads implied: the gap is the premium. Deals
broke 13.7\% of the time when the market fell over their life (117 deals) and 4.7\% when it rose. The portfolio's beta is 0.20 in the ten months when the
market fell by more than 4\%, and 0.09 in the other 86: the put-like profile, weaker than in the historical data because the synthetic market had few
severe falls. Two counterfactual columns separate the pieces. If breaks did not depend on the market, the same spreads would earn 4.30\%: the clustering
of breaks in downturns costs 1.5 points a year. If spreads priced only the flat-market break rate, the portfolio would lose 0.47\% a year: the premium in
the spread is exactly the price of the crash risk.

\begin{omfigure}
\begin{tikzpicture}
\begin{axis}[width=8.5cm,height=6.0cm,xlabel={\small market return in the month (\%)},ylabel={\small deal portfolio return (\%)},
  tick label style={font=\footnotesize},xmin=-16,xmax=12,ymin=-4,ymax=3.5,grid=major,grid style={gray!25},table/col sep=comma]
\addplot[omDef,only marks,mark=*,mark size=1.4pt] table[x=market_pct,y=deals_pct]{figdata/strategies-1/11-merger-arbitrage/months.csv};
\addplot[black!40,dashed] coordinates {(-4,-4) (-4,3)};
\end{axis}
\end{tikzpicture}

\omcaption{Monthly returns of the synthetic deal portfolio against the market's, years 3 to 10; the dashed line marks a 4\% market fall. Data:
\texttt{s1\_mergerarb.run}.}\label{fig:s1:merger-arbitrage:months}
\end{omfigure}

The worst month for the portfolio ($-3.06\%$) was a month in which the market rose 2.8\%: with twenty deals, two or three breaking together is enough, and
deal-specific breaks are the other half of the risk. A merger-arbitrage book diversifies that half by holding many deals; it cannot diversify the half that
comes from the market.

\section{Stock deals and collars}

\begin{definition}[Stock-deal collar]\label{def:s1:merger-arbitrage:collar}
A \emph{stock-deal collar}\index{stock-deal collar} fixes the exchange ratio of a stock deal while the acquirer's price stays inside a band, and fixes the
value of the consideration outside it (the ratio then adjusts), protecting the target's shareholders against large falls in the acquirer's price.
\end{definition}

In a stock deal the target's holders receive a number of acquirer shares; the arbitrageur buys the target and sells short that number of acquirer shares,
so that the spread no longer depends on the acquirer's price. A collar changes the hedge. With a ratio of 0.5 and a band of \$40 to \$50, an acquirer price of
\$42 gives consideration of \$21, hedged by shorting 0.5 acquirer shares; below \$40, at \$36, the consideration is fixed at \$20 and the hedge is none, since
the value no longer moves with the acquirer. The hedge ratio is the collar's delta (\texttt{collar\_hedge}), and near the band's edges it jumps: collar
deals need rebalancing, and their spreads carry the gamma of that option.

\section{Regulatory and financing risk}

Deals break for reasons the market's level does not capture. Antitrust review (the waiting periods in the dated box, and a Second Request that can extend
them for months) is the best-known; financing that disappears, a shareholder vote that fails, a competing bid (which raises the price instead) and material
adverse changes in the target's business are the others. Each has its own probability and its own \omterm{def:s1:merger-arbitrage:spread}{break price}, and a deal's spread is a mix of them. The
synthetic model folds them into one probability with a market-dependent part; a real desk estimates them deal by deal, reads the merger agreement, and sizes
positions by what each deal can lose, not by its spread.

\section{Portfolio of deals}

\omterm{def:s1:merger-arbitrage:arb}{Merger arbitrage} is limited by capital as much as by opportunity. Baker and Savasoglu studied the limits to that capital in
mergers and acquisitions; Mitchell and Pulvino's put-like payoff is one reason it is limited. A deal book's risk budget has three lines: deal-specific
breaks, which diversify across twenty or more deals; market-dependent breaks, which do not; and the market exposure of the \omterm{def:s1:merger-arbitrage:spread}{break prices} themselves, which
shows up as the book's beta in falls. The last two are the insurance sold. The honest benchmark for the book is a short put on the index, not cash.

\section{Strategy files}

\begin{strategyfile}{Cash-deal spread}\label{strat:s1:merger-arbitrage:cash}
\sfield{Who pays you, and why} Target shareholders who sell after the announcement rather than wait and bear the break risk.
\sfield{Instruments and venues} The target's shares.
\sfield{Signal} The annualised spread against an estimate of the break probability and \omterm{def:s1:merger-arbitrage:spread}{break price}.
\sfield{Sizing and execution} Size by the loss if the deal breaks, not by the spread; enter after the announcement.
\sfield{Costs} Low turnover; financing of the position.
\sfield{How it dies} Breaks, especially in downturns; regulatory surprises.
\sfield{Horizon, capacity, infrastructure} Months; capacity limited by deal size; legal and regulatory analysis.
\sfield{Backtest honestly} Deals as announced, including those that broke; \omterm{def:s1:merger-arbitrage:spread}{break prices} estimated before the outcome.
\sfield{Sources} Mitchell and Pulvino (2001): 4\,750 mergers, 1963--1998, excess returns of about 4\% a year after costs, put-like payoff.
\end{strategyfile}

\begin{strategyfile}{Stock-for-stock hedged spread}\label{strat:s1:merger-arbitrage:stock}
\sfield{Who pays you, and why} As for cash deals, with the acquirer's price risk hedged.
\sfield{Instruments and venues} The target, long; the acquirer, short in the exchange ratio.
\sfield{Signal} The spread between the target's price and the ratio times the acquirer's.
\sfield{Sizing and execution} Keep the short at the exchange ratio; adjust for dividends.
\sfield{Costs} Borrow on the acquirer, sometimes expensive when many arbitrageurs short it.
\sfield{How it dies} Breaks; a squeeze in the acquirer if many are short.
\sfield{Horizon, capacity, infrastructure} Months; borrow availability.
\sfield{Backtest honestly} Borrow fees and recalls on the acquirer; the acquirer's move on a break.
\sfield{Sources} Mitchell and Pulvino (2001).
\end{strategyfile}

\begin{strategyfile}{Collar deals}\label{strat:s1:merger-arbitrage:collar}
\sfield{Who pays you, and why} As for stock deals, with an option embedded in the consideration.
\sfield{Instruments and venues} The target, the acquirer, and possibly listed options to hedge the collar's gamma.
\sfield{Signal} The spread against the collar's value.
\sfield{Sizing and execution} A hedge ratio equal to the collar's delta, rebalanced near the band's edges.
\sfield{Costs} Rebalancing of the hedge.
\sfield{How it dies} Hedging errors near the band's edges; breaks.
\sfield{Horizon, capacity, infrastructure} Months; an option model of the collar.
\sfield{Backtest honestly} The collar's terms as written, including averaging periods.
\sfield{Sources} This chapter's collar arithmetic; no performance record verified.
\end{strategyfile}

\begin{strategyfile}{Diversified deal portfolio}\label{strat:s1:merger-arbitrage:portfolio}
\sfield{Who pays you, and why} The market, for bearing concentrated break risk in downturns.
\sfield{Instruments and venues} Twenty or more announced deals.
\sfield{Signal} Each deal's spread against its estimated risk.
\sfield{Sizing and execution} Equal risk per deal; limits on deals in the same industry or with the same regulator.
\sfield{Costs} Low.
\sfield{How it dies} Market crashes, when breaks cluster and \omterm{def:s1:merger-arbitrage:spread}{break prices} fall together.
\sfield{Horizon, capacity, infrastructure} Continuous; capacity set by the deal flow.
\sfield{Backtest honestly} All deals, point in time; compare with a short index put.
\sfield{Sources} Mitchell and Pulvino (2001); this chapter's simulation (2.76\% a year over the rate, beta 0.20 in falling months and 0.09 otherwise).
\end{strategyfile}

\section{Tutorial: selling insurance}\label{tut:s1:merger-arbitrage:tut}

\begin{tutorial}
\textbf{Goal.} Price \omterm{def:s1:merger-arbitrage:spread}{deal spreads}, simulate a portfolio of deals whose breaks depend on the market, and measure its payoff. \textbf{End state:} the table
and \cref{fig:s1:merger-arbitrage:months}.
\begin{tutsteps}
\item \textbf{Pricing}: implied probability and the collar.
\omcode{code/firm/mergerarb/firm_mergerarb.py}{36}{48}{Implied completion probability and collar terms.}\label{lst:s1:merger-arbitrage:pricing}
\item \textbf{Deals}: announcement, price path, break or completion.
\omcode{code/firm/mergerarb/firm_mergerarb.py}{65}{88}{A deal's life, with market-dependent breaks.}\label{lst:s1:merger-arbitrage:deals}
\item \textbf{Run} \texttt{run()} and its counterfactuals, \texttt{break\_rates()} and \texttt{fig\_mergerarb.py}.
\end{tutsteps}
\textbf{What to change next.} Add stock deals with collars and hedge them; let spreads widen after market falls (the market reprices $q$); cap the book's
exposure to deals in one industry.
\end{tutorial}

\section{Build: merger arbitrage}\label{bld:s1:merger-arbitrage:mergerarb}

\begin{build}
\bfield{Purpose} Deal-spread pricing and a portfolio of deals with market-dependent breaks.
\bfield{Interface} \texttt{spread(price, offer)}, \texttt{annualised(spread, days)}, \texttt{implied\_probability(price, offer, break\_price, rate, days)},
\texttt{collar\_value(acquirer, ratio, low, high)}, \texttt{collar\_hedge(acquirer, ratio, low, high)}, \texttt{DealConfig}, \texttt{simulate\_deals(mkt, cfg, rng)},
\texttt{portfolio(deals, T)}.
\bfield{Rules} Deals entered after their announcement; breaks decided by information at the deal's end; returns in excess of the discount rate reported.
\bfield{Acceptance tests} \texttt{code/firm/mergerarb/tests/}: spread, implied probability and annualisation by hand; collar value and hedge; deals in a flat
market break at the base rate and the portfolio earns a positive return.
\bfield{Stretch} Stock deals with acquirer hedges; competing bids; spreads that widen when the market falls.
\end{build}

\begin{omsources}
\item M.~Mitchell and T.~Pulvino, ``Characteristics of risk and return in risk arbitrage'', \emph{Journal of Finance} 56(6), 2001.
\item M.~Baker and S.~Savasoglu, ``Limited arbitrage in mergers and acquisitions'', \emph{Journal of Financial Economics} 64(1), 2002.
\item US Federal Trade Commission, Premerger Notification and the Merger Review Process.
\end{omsources}

\section{Exercises}

\begin{exercise}[$\star$]\label{exo:s1:merger-arbitrage:1}
A cash offer of \$40 trades at \$38.60 with a \omterm{def:s1:merger-arbitrage:spread}{break price} of \$30. Compute the spread and the implied completion probability.
\end{exercise}

\begin{exercise}[$\star$]\label{exo:s1:merger-arbitrage:2}
The deal is expected to close in 84 trading days. Annualise the spread.
\end{exercise}

\begin{exercise}[$\star$]\label{exo:s1:merger-arbitrage:3}
A stock deal pays 0.5 acquirer shares with a collar from \$40 to \$50. What is the consideration at acquirer prices of \$42 and \$36, and how many acquirer
shares do you short per target share in each case?
\end{exercise}

\begin{exercise}[$\star\star$]\label{exo:s1:merger-arbitrage:4}
Show that buying the \$40 deal at \$38.60 has zero expected profit at the implied probability, and state what the arbitrageur must believe to trade.
\end{exercise}

\begin{exercise}[$\star\star$]\label{exo:s1:merger-arbitrage:5}
Why does a deal portfolio's beta rise in falling markets even if its deals are unrelated to each other?
\end{exercise}

\begin{exercise}[$\star\star$]\label{exo:s1:merger-arbitrage:6}
What does a Second Request do to a deal's spread, and why?
\end{exercise}

\begin{exercise}[$\star\star\star$]\label{exo:s1:merger-arbitrage:7}
\emph{Coding.} Run \texttt{run(0.0, 0.12, 0.06)} and \texttt{run(1.5, 0.06, 0.06)}. Explain what each counterfactual isolates.
\end{exercise}

\begin{exercise}[$\star\star\star$]\label{exo:s1:merger-arbitrage:8}
\emph{Find the flaw.} ``\omterm{def:s1:merger-arbitrage:arb}{Merger arbitrage} has a Sharpe ratio near one and a beta near zero; it is market neutral.''
\end{exercise}

\section{Problem: Selling Insurance}

\begin{problem}[{Weekend problem --- the premium and the claims}]\label{pb:s1:merger-arbitrage:1}
The chapter's synthetic deals and the public record.

\medskip\noindent\textbf{Part I --- Pricing.}
\begin{enumerate}
\item Define \omterm{def:s1:merger-arbitrage:arb}{merger arbitrage}, a \omterm{def:s1:merger-arbitrage:arb}{tender offer}, the \omterm{def:s1:merger-arbitrage:spread}{deal spread} and the \omterm{def:s1:merger-arbitrage:spread}{break price}.
\item Compute the implied probability of the opening deal.
\item Why are spreads annualised?
\item Describe the synthetic deals.
\end{enumerate}

\medskip\noindent\textbf{Part II --- The portfolio.}
\begin{enumerate}[resume]
\item Give its return, excess return, volatility and Sharpe ratio.
\item Compare the implied and realised completion rates.
\item How do break rates depend on the market?
\item Give the betas in falling and other months.
\end{enumerate}

\medskip\noindent\textbf{Part III --- The pieces.}
\begin{enumerate}[resume]
\item What do the two counterfactuals show?
\item Describe the worst month and what it says about deal-specific risk.
\item How are stock deals and collars hedged?
\item What are the non-market reasons deals break?
\end{enumerate}

\medskip\noindent\textbf{Part IV --- The verdict.}
\begin{enumerate}[resume]
\item State the \emph{named result}: the portfolio's beta in down and up markets and its implied completion probability against realised completion.
\item What did Mitchell and Pulvino find?
\item What is the right benchmark for a deal book?
\item How would you size individual deals?
\item What does the HSR waiting period mean for a deal's timeline?
\item Which risk can a deal book diversify, and which not?
\item When would you add to the book: after a market fall or after a rise?
\item In one sentence: what does a merger arbitrageur sell?
\end{enumerate}
\end{problem}

\section{Interview questions}

\begin{interviewq}[$\star$ \iqroles{trader, researcher}]\label{iq:s1:merger-arbitrage:1}
What is \omterm{def:s1:merger-arbitrage:arb}{merger arbitrage}, and where does its return come from?
\end{interviewq}

\begin{interviewq}[$\star\star$ \iqroles{trader}]\label{iq:s1:merger-arbitrage:2}
How do you hedge a stock-for-stock deal with a collar?
\end{interviewq}

\begin{interviewq}[$\star\star$ \iqroles{risk}]\label{iq:s1:merger-arbitrage:3}
Why is a merger-arbitrage book compared to a short put?
\end{interviewq}

\begin{interviewq}[$\star\star$ \iqroles{researcher}]\label{iq:s1:merger-arbitrage:4}
How would you estimate a deal's break probability and \omterm{def:s1:merger-arbitrage:spread}{break price}?
\end{interviewq}

\begin{interviewq}[$\star\star$ \iqroles{trader}]\label{iq:s1:merger-arbitrage:5}
A deal's spread doubles overnight with no news. What might have happened?
\end{interviewq}

\begin{interviewq}[$\star\star\star$ \iqroles{researcher}]\label{iq:s1:merger-arbitrage:6}
Derive the implied completion probability with discounting at rate $r$ over $\tau$ years, and show how the arbitrageur's expected excess return depends on
the gap between the implied and the true probability.
\end{interviewq}
